% SPDX-License-Identifier: Apache-2.0
\section{A Wait Under If}
\label{sec:x-frame}

Section~\ref{sec:x-seq} numbered a loop's waits into states and
Section~\ref{sec:x-serial} unrolled a counted step. The step left is
the one taken only sometimes: a wait under \code{if}. Its state is
made like any other, and the state before the \code{if} gets a way
around it. When the condition holds at that state's edge the
register goes to the wait inside; when it fails, the register goes
past the \code{if}, to whatever wait comes next. So the register's
next value is a choice on the condition rather than a number, and
that choice is the whole of the control flow the \code{if} asked
for.

The statements around the \code{if} keep their meaning by being
placed where they happen. What sits in an arm before its wait
happens in the state before the \code{if}, under the arm's
condition and the failure of the arms above it; what sits in an arm
after its wait happens in the state the wait made, under nothing,
since being there says the condition held. What follows the
\code{if} happens in every state the \code{if} can end in: the last
state of each arm, and the state before the \code{if} when there is
no \code{else}, under the failure of every arm. It is lowered once,
under the union of those states' conditions, rather than once per
state, so a tail shared by many states costs one copy of its logic
and one wire per \code{let} in it. An \code{if} inside
an \code{if} nests the same way, and a \code{for} that waits inside
an arm is unrolled there first. A \code{let} still does not cross a
wait, and an output port still holds what the last state set, as
Section~\ref{sec:x-seq} says.

The framer below sends a byte as a serial frame: a start bit, eight
data bits low bit first, a parity bit when \code{parity} asks for
one at the last data bit, and a stop bit. Twelve states, and the
parity state is the one under \code{if}; the netlist shows the
choice, \code{parity ? 10 : 11}, where the state of the last data
bit picks its successor. The build checks the netlist against the
run under nvc and Verilator, with one frame of each length.

\begin{figure}[htbp]
\centering
\input{frame_timing.tex}
\caption{The framer: a start bit, the eight bits of the byte, and a
stop bit; the second frame has a parity bit before its stop bit,
since \code{parity} is high by then.}
\label{fig:x-frame}
\end{figure}

\lstinputlisting[caption={\code{ex\_frame.rs}. Run with \code{bazel run //lib/examples:ex\_frame}.},label={lst:x-frame}]{lib/examples/ex_frame.rs}

\lstinputlisting[style=ascii,caption={What \code{ex\_frame} printed when the build ran it: the two frames bit by bit, then the Verilog, with the choice of the next state in it.},label={lst:x-frame-out}]{out_frame.txt}
